\documentclass[12pt, a4paper]{article}

\usepackage[T2A]{fontenc}
\usepackage[utf8]{inputenc}
\usepackage[bulgarian]{babel}
\usepackage{geometry}
\geometry{a4paper, margin=1in}
\usepackage{tabularx}
\usepackage{ragged2e}
\usepackage{hyperref}

\title{Задание 5: Курсов проект (Система, изходни кодове, отчет)}
\author{}
\date{}

\begin{document}

\maketitle

\noindent\textbf{Дисциплина:} Проектиране на човеко-машинен интерфейс 2025-2026

\bigskip

\noindent\textbf{Участници в проекта:}
\smallskip

\noindent
\begin{tabularx}{\textwidth}{|l|X|l|X|l|}
\hline
\textbf{№} & \textbf{Име и фамилия} & \textbf{Факултетен №} & \textbf{Специалност} & \textbf{Курс} \\ \hline
1 & Виктор Койчев & 3MI0600164 & Софтуерно инженерство & 3-ти \\ \hline
2 & Емилиан Янев & 0MI0600273 & Софтуерно инженерство & 3-ти \\ \hline
3 & Жан Георгиев & 8MI0600398 & Софтуерно инженерство & 3-ти \\ \hline
4 & Даниел Николов & 0MI0600410 & Софтуерно инженерство & 3-ти \\ \hline
\end{tabularx}

\bigskip

\noindent\textbf{Име на проекта:}\\
Система за управление и мониторинг на роботизиран склад за фармацевтични продукти (PharmaBot Logistics)

\bigskip
\rule{\textwidth}{0.4pt}

\section*{1. Описание на системата}
Курсовият проект \textit{PharmaBot Logistics} представлява интерактивен (Low-Fidelity) уеб прототип, демонстриращ управлението на флотилия от автономни складови роботи. Системата е проектирана с фокус върху потребителското преживяване (UX/UI) на две основни групи: мениджъри/оператори в склада и шофьори на товарните рампи. Прототипът е реализиран чрез HTML, Tailwind CSS и Vanilla JavaScript, което позволява директно тестване на основните потребителски взаимодействия в браузърна среда.

\section*{2. Достъп до проекта}

\begin{itemize}
    \item \textbf{Работеща система (URL):} \url{https://veector40.github.io/PCHMI/index.html}
    \item \textbf{Изходни кодове (GitHub Repository):} \url{https://github.com/Veector40/PCHMI}
\end{itemize}

\section*{3. Ръководство за тестване (Сценарии)}

Моля, отворете посочения URL адрес и следвайте стъпките по-долу, за да тествате интегрираните функционалности:

\subsection*{Сценарий 1: Работа на складовия оператор (Десктоп)}
\begin{enumerate}
    \item На началния екран въведете поне 3 символа в полето за 2FA код (за симулация на сигурност).
    \item Кликнете върху бутона \textbf{„Вход като Оператор / Мениджър“}. Ще бъдете пренасочени към контролния панел (\texttt{dashboard.html}).
    \item Разгледайте \textbf{2D Картата}. Задръжте курсора на мишката върху движещия се (зелен) или зареждащия се (жълт) робот, за да видите детайлен Tooltip с информация.
    \item Тествайте функцията за \textbf{Geo-fencing}: Кликнете бутона \textbf{„Режим: Очертай Зона“} горе вдясно. След това кликнете някъде върху картата, за да създадете визуална зона на разлив.
    \item Тествайте системата за аларми: Натиснете бутона \textbf{„Тест: Температурна Аларма“}, за да видите модалния прозорец за критичен инцидент.
    \item Използвайте бутона \textbf{„Изход от системата“} долу вляво и потвърдете, за да се върнете на екрана за вход.
\end{enumerate}

\subsection*{Сценарий 2: Приемане на пратка (Мобилен изглед)}
\begin{enumerate}
    \item На екрана за вход въведете 2FA код и изберете \textbf{„Вход като Шофьор (Рампа)“}.
    \item Ще се зареди мобилен изглед (\texttt{mobile-handoff.html}). Ако симулирате работа на открито (на пряка слънчева светлина), тествайте бутона \textbf{„Контраст“} (горе вдясно) за преминаване в специален четим режим.
    \item Забележете, че Стъпка 2 е заключена. Кликнете върху \textbf{„Сканирай QR код“} (Стъпка 1). 
    \item След симулираното сканиране, системата ще отключи \textbf{„Снимай палета“} (Стъпка 2). Кликнете го.
    \item Едва след това финалният бутон \textbf{„Потвърди Експедиция“} ще стане активен. Кликването му завършва процеса и ви връща на началния екран.
\end{enumerate}

\end{document}